% =====================================================================
%  plan-de-proyecto.tex  --  Plan de Proyecto de LEIP
%  ARCHIVO UNICO: no necesita ningun otro archivo.
%  Compilar en Overleaf con pdfLaTeX (dos pasadas por el indice).
% =====================================================================
\documentclass[11pt,a4paper]{article}

% ---------------------------------------------------------------------
%  DATOS DEL PROYECTO
% ---------------------------------------------------------------------
\newcommand{\proyectoNombre}{LEIP}
\newcommand{\proyectoNombreLargo}{LatAm Economic Intelligence Platform}
\newcommand{\proyectoUniversidad}{Universidad Cat\'olica del Uruguay}
\newcommand{\proyectoFacultad}{Facultad de Ingenier\'ia y Tecnolog\'ias}
\newcommand{\proyectoCarrera}{Ingenier\'ia en Inform\'atica}
\newcommand{\docVersion}{1.0}
\newcommand{\docTitulo}{Plan de Proyecto}
\newcommand{\docFecha}{Agosto 2026}

% ---------------------------------------------------------------------
%  PREAMBULO
% ---------------------------------------------------------------------
\usepackage[utf8]{inputenc}
\usepackage[T1]{fontenc}
\IfFileExists{lmodern.sty}{\usepackage{lmodern}\usepackage{microtype}}{}

\IfFileExists{spanish.ldf}{%
  \usepackage[spanish,es-tabla,es-noquoting]{babel}%
}{%
  \usepackage[english]{babel}%
  \renewcommand{\contentsname}{\'Indice}%
  \renewcommand{\listfigurename}{\'Indice de figuras}%
  \renewcommand{\listtablename}{\'Indice de tablas}%
  \renewcommand{\tablename}{Tabla}%
  \renewcommand{\figurename}{Figura}%
}

\usepackage[a4paper,margin=2.5cm,headheight=15pt]{geometry}
\usepackage{fancyhdr}
\usepackage{parskip}
\usepackage{booktabs}
\usepackage{tabularx}
\usepackage{longtable}
\usepackage{array}
\usepackage{enumitem}
\setlist{noitemsep,topsep=4pt}
\usepackage{graphicx}
\usepackage{xcolor}
\usepackage[hidelinks]{hyperref}
\usepackage{url}
\urlstyle{same}

\definecolor{leipAzul}{HTML}{1F4E79}
\definecolor{leipGris}{HTML}{5A5A5A}

\pagestyle{fancy}
\fancyhf{}
\fancyhead[L]{\small\color{leipGris}\proyectoNombre{} --- \docTitulo}
\fancyhead[R]{\small\color{leipGris} v\docVersion}
\fancyfoot[C]{\small\thepage}
\renewcommand{\headrulewidth}{0.4pt}

\usepackage{titlesec}
\titleformat{\section}{\Large\bfseries\color{leipAzul}}{\thesection}{0.6em}{}
\titleformat{\subsection}{\large\bfseries\color{leipAzul}}{\thesubsection}{0.6em}{}
\titleformat{\subsubsection}{\normalsize\bfseries\color{leipGris}}{\thesubsubsection}{0.6em}{}

% --- Portada -----------------------------------------------------------
\newcommand{\portada}[1]{%
  \begin{titlepage}
  \centering
  \vspace*{2cm}
  {\Large \proyectoUniversidad\par}
  {\large \proyectoFacultad\par}
  {\large \proyectoCarrera\par}
  \vspace{3cm}
  {\Huge\bfseries\color{leipAzul} #1\par}
  \vspace{0.8cm}
  {\Large \proyectoNombre{} --- \proyectoNombreLargo\par}
  \vspace{2cm}
  {\large Versi\'on \docVersion\par}
  {\large \docFecha\par}
  \vfill
  {\large\bfseries Integrantes\par}
  \vspace{0.3cm}
  \begin{tabular}{c}
    Sebastian Forische \\
    Armando Hernandez de la Vega \\
    Luana Acu\~na \\
    Bruno Sebastian Olivera Viera Gomez \\
  \end{tabular}
  \vspace{1.5cm}
  \end{titlepage}
}

% --- Historial de revisiones -------------------------------------------
\newenvironment{historial}{%
  \section*{Historial de revisiones}
  \addcontentsline{toc}{section}{Historial de revisiones}
  \begin{tabular}{@{}l l p{0.52\textwidth} l@{}}
  \toprule
  \textbf{Versi\'on} & \textbf{Fecha} & \textbf{Descripci\'on} & \textbf{Autor} \\
  \midrule
}{%
  \bottomrule
  \end{tabular}
  \bigskip
}

% =====================================================================
%  ENTORNO PARA DESCRIBIR UN SPRINT
% ---------------------------------------------------------------------
%  Uso:
%
%  \begin{sprint}{1}{Del 01/09 al 21/09}
%    \sobjetivo{Objetivo del sprint}
%    \salcance{Que casos de uso o requerimientos entran}
%    \sentregable{Que se entrega al cierre}
%    \sriesgos{Que riesgos de la Lista de Riesgos se mitigan}
%  \end{sprint}
% =====================================================================

\newcommand{\scampo}[2]{%
  \par\noindent
  \begin{tabularx}{\textwidth}{@{}>{\bfseries}p{3.4cm} X@{}}
    #1 & #2
  \end{tabularx}\par\smallskip
}

\newcommand{\sobjetivo}[1]{\scampo{Objetivo}{#1}}
\newcommand{\salcance}[1]{\scampo{Alcance}{#1}}
\newcommand{\sentregable}[1]{\scampo{Entregable}{#1}}
\newcommand{\sriesgos}[1]{\scampo{Riesgos que mitiga}{#1}}

\newenvironment{sprint}[2]{%
  \subsection*{\color{leipAzul}Sprint #1\quad\normalsize\color{leipGris}#2}
  \addcontentsline{toc}{subsection}{Sprint #1}
  \vspace{-0.4em}
  {\color{leipGris}\hrule height 0.5pt}
  \smallskip
}{%
  \bigskip
}

% --- Figura con fallback si la imagen todavia no existe ----------------
\newcommand{\figuraLEIP}[4]{%
  \begin{figure}[htbp]
  \centering
  \IfFileExists{#1}{\includegraphics[width=#4\textwidth]{#1}}{%
    \fbox{\parbox[c][3.5cm][c]{0.8\textwidth}{\centering\color{leipGris}%
      \small\texttt{#1}}}%
  }
  \caption{#2}
  \label{#3}
  \end{figure}
}

\begin{document}

\portada{Plan de Proyecto}

\tableofcontents
\bigskip

\begin{historial}
1.0 & 26/06/2026 & Contenido inicial derivado del anteproyecto & Equipo \\
\end{historial}
\clearpage

% =====================================================================
\section{Introducci\'on}

\subsection{Objetivo del documento}

\subsection{Alcance}

% =====================================================================
\section{Marco de trabajo}

El equipo trabaja bajo una metodolog\'ia \'agil (Scrum), con iteraciones cortas
(\emph{sprints}) para la planificaci\'on, desarrollo, validaci\'on y revisi\'on
continua del proyecto, manteniendo tambi\'en un tablero de \emph{backlog}.

\subsection{Duraci\'on de los sprints}

\subsection{Ceremonias}

\subsection{Definici\'on de terminado}

% =====================================================================
\section{Equipo y responsabilidades}

\begin{table}[htbp]
\centering
\caption{Responsables por \'area}
\label{tab:equipo}
\begin{tabularx}{\textwidth}{@{}l X@{}}
\toprule
\textbf{\'Area} & \textbf{Encargado} \\
\midrule
Gesti\'on del proyecto          & Bruno Olivera, Luana Acu\~na \\
An\'alisis de requerimientos    & Armando Hernandez, Bruno Olivera \\
Dise\~no de la interfaz         & Bruno Olivera, Sebastian Forische \\
Desarrollo frontend             & Luana Acu\~na, Bruno Olivera \\
Desarrollo backend              & Sebastian Forische, Luana Acu\~na, Bruno Olivera \\
Infraestructura                 & Sebastian Forische, Luana Acu\~na \\
QA                              & Armando Hernandez, Luana Acu\~na \\
\emph{Database Administrator}   & Sebastian Forische \\
Procesamiento de datos con LLM  & Armando Hernandez, Sebastian Forische \\
\emph{Product Owner}            & Armando Hernandez \\
\bottomrule
\end{tabularx}
\end{table}

\subsection{Roles de Scrum}

% =====================================================================
\section{Product backlog}

\begin{tabularx}{\textwidth}{@{}l X c c@{}}
\toprule
\textbf{ID} & \textbf{Historia de usuario} & \textbf{Requerimientos} & \textbf{Prioridad} \\
\midrule
 &  &  &  \\
 &  &  &  \\
 &  &  &  \\
\bottomrule
\end{tabularx}

\subsection{Criterios de priorizaci\'on}

% =====================================================================
\section{Planificaci\'on de sprints}

% ---------------------------------------------------------------------
%  PLANTILLA: copiá este bloque para cada sprint.
% ---------------------------------------------------------------------
\begin{sprint}{0}{Fechas}

\sobjetivo{}
\salcance{}
\sentregable{}
\sriesgos{}

\end{sprint}

% =====================================================================
\section{Cronograma}

\figuraLEIP{img/cronograma.png}{Cronograma del proyecto}{fig:cronograma}{0.95}

% =====================================================================
\section{Estimaci\'on}

\subsection{M\'etodo de estimaci\'on}

\subsection{Velocidad del equipo}

% =====================================================================
\section{Gesti\'on de cambios}

% =====================================================================
\section{Comunicaci\'on y reuniones}

% =====================================================================
\section{Herramientas}

\begin{tabularx}{\textwidth}{@{}l X@{}}
\toprule
\textbf{Herramienta} & \textbf{Uso} \\
\midrule
 &  \\
 &  \\
 &  \\
\bottomrule
\end{tabularx}

% =====================================================================
\section{Criterio de aceptaci\'on del proyecto}

% =====================================================================
\section{Gesti\'on de riesgos}

\end{document}
